% References and claim scope: experiments/limitations_sources.md.
\renewcommand{\slideid}{L3}
\section{Why classical solvers remain essential}

\renewcommand{\slideid}{L3-C01}
\begin{frame}[c]{Why classical numerical methods are still used}
\small
\guidingquestion{Can we reach a requested error tolerance at a predictable cost?}
\begin{slideitems}
\item \textbf{Controlled refinement:} mesh size, polynomial degree, and time step
give systematic ways to improve a discretization.
\item \textbf{Efficient algebra:} sparse matrices, preconditioners, and multigrid
exploit the structure of many PDEs.
\item \textbf{Physical structure:} appropriate finite-volume, mixed, or compatible
schemes can enforce conservation or other identities by construction.
\item \textbf{Established practice:} error estimation, complex geometry handling,
and extensively tested implementations.
\end{slideitems}
\begin{block}{A published benchmark}
In the PDE benchmarks of Grossmann et al., PINNs did not beat FEM in
solution time and accuracy on the problems tested.
\end{block}
\labnote{\href{https://arxiv.org/abs/2302.04107}{Grossmann et al., IMA J. Applied Mathematics (2024).}}
\end{frame}

\renewcommand{\slideid}{L3-C02}
\begin{frame}[c]{A classical route from approximation to accuracy}
\small
For conforming FEM ($V_h\subset V$) with bounded, coercive $a$,
C\'ea's lemma gives
\[
 \norm{u-u_h}_V\le \frac{M}{\alpha}
       \inf_{v_h\in V_h}\norm{u-v_h}_V.
\]
Here $M$ is the continuity constant and $\alpha>0$ the coercivity constant;
$u_h$ is the exactly solved Galerkin approximation.
\begin{slideitems}
\item With sufficient regularity and suitable meshes, approximation estimates
turn this into an explicit rate in $h$ and polynomial degree.
\item Reliable a posteriori estimators are available for many problem classes
and can guide local refinement.
\end{slideitems}
\begin{block}{A finite-tolerance solve still needs an error budget}
\[
\norm{u-\widetilde u_h}_V
\le \underbrace{\norm{u-u_h}_V}_{\text{discretization}}
  + \underbrace{\norm{u_h-\widetilde u_h}_V}_{\text{algebraic solve}}.
\]
Nonlinear, singular, or unstable problems still require appropriate analysis.
\end{block}
\end{frame}

\renewcommand{\slideid}{L3-C03}
\begin{frame}[c]{What a PINN error bound requires}
\small
\begin{block}{An error estimate under two assumptions}
If PDE stability gives $\norm{u-u_\theta}_X\le C_{\rm stab}\sqrt{\Loss(\theta)}$
and sampling gives $\Loss(\theta)\le\widehat\Loss(\theta)+\varepsilon_Q$, then
\[
 \norm{u-u_\theta}_X\le
 C_{\rm stab}\sqrt{\widehat\Loss(\theta)+\varepsilon_Q}.
\]
The loss must include the right residual and boundary norms for this PDE.
\end{block}
\begin{slideitems}
\item \textbf{Representation:} a good network may exist; training need not find it.
\item \textbf{Sampling:} a useful $\varepsilon_Q$ requires regularity and coverage,
with control valid for the network selected by training.
\item \textbf{Optimization:} a bound assuming small training error does not
prove Adam or L-BFGS reaches it at an affordable cost.
\end{slideitems}
\labnote{\href{https://arxiv.org/abs/2006.16144}{Mishra--Molinaro (2023)};
\href{https://doi.org/10.1017/S0962492923000089}{De Ryck--Mishra, Acta Numerica (2024)}.
Norms, rates, and constants depend on the PDE and formulation.}
\end{frame}

\renewcommand{\slideid}{L3-C05}
\begin{frame}[c]{Small training loss can still hide a wrong solution}
\small
\renewcommand{\arraystretch}{1.35}
\begin{tabular}{@{}>{\raggedright\arraybackslash}p{0.30\textwidth}>{\raggedright\arraybackslash}p{0.65\textwidth}@{}}
\toprule
Difficulty & What can go wrong\\
\midrule
Finite collocation & Residual peaks or thin layers fall between training points.\\
Multiple scales & Oscillations and sharp features can learn much more slowly
than the smooth background.\\
Time dependence & Small early-time errors and an unbalanced space--time loss
can spoil long-time predictions.\\
Physical constraints & A generic pointwise loss does not exactly impose cell
conservation, positivity, or an entropy inequality.\\
\bottomrule
\end{tabular}
\begin{block}{Adding data does not remove these issues}
Noisy or sparse observations introduce their own weighting and
identifiability questions, as the inverse experiment demonstrated.
\end{block}
\labnote{\href{https://arxiv.org/abs/2109.01050}{Krishnapriyan et al., NeurIPS (2021)};
\href{https://arxiv.org/abs/2007.14527}{Wang--Yu--Perdikaris, JCP (2022)}.
These are possible failure modes, not unavoidable outcomes.}
\end{frame}

\renewcommand{\slideid}{L3-C09}
\begin{frame}[c]{What helps PINNs, and what still needs checking}
\small
\renewcommand{\arraystretch}{1.4}
\begin{tabular}{@{}>{\raggedright\arraybackslash}p{0.41\textwidth}>{\raggedright\arraybackslash}p{0.54\textwidth}@{}}
\toprule
Possible improvement & Remaining check\\
\midrule
Hard boundary constraints & Interior PDE accuracy and ansatz suitability.\\
Adaptive collocation & Independent residual evaluation and correct sampling weights.\\
Feature scaling, domain splitting & Resolution of every relevant physical scale.\\
Better conditioning and optimizers & Actual solution error; convergence is still problem-dependent.\\
Weak or variational formulations & Stability, quadrature, and the chosen test space.\\
\bottomrule
\end{tabular}
\begin{block}{Check whether the change helps}
Measure the solution error before and after each change. Improvements
depend on the PDE, sampling, and training setup.
\end{block}
\labnote{\href{https://doi.org/10.1017/S0962492923000089}{De Ryck--Mishra (2024)};
\href{https://proceedings.mlr.press/v235/rathore24a.html}{Rathore et al. (2024)}.}
\end{frame}

\renewcommand{\slideid}{L3-C11}
\begin{frame}[c]{Where neural PDE solvers can have an advantage}
\small
\begin{slideitems}
\setlength\itemsep{0.8em}
\item \textbf{High-dimensional settings.} Neural representations and stochastic
sampling avoid the exponential size of tensor grids. For PDEs with exploitable
structure, this can make otherwise infeasible computations practical.
\item \textbf{Inverse problems with sparse or poorly distributed data.}
Jointly fit the state and unknown coefficients, using the PDE to constrain
regions with few measurements. This can avoid repeated full forward solves;
identifiability still requires informative observations.
\item \textbf{Complex geometry.} Meshfree formulations use interior and boundary
samples, reducing volume-meshing and remeshing effort for intricate or changing
domains. Accurate boundary representation and enforcement remain essential.
\end{slideitems}
\par\medskip
Compare total cost at matched accuracy, including training and geometry preparation.
\par\medskip
\labnote{Examples: \href{https://arxiv.org/abs/1707.02568}{Han--Jentzen--E (2018)};
\href{https://doi.org/10.1126/science.aaw4741}{Raissi et al. (2020)};
\href{https://arxiv.org/abs/2104.08426}{Sukumar--Srivastava (2022)}.}
\end{frame}

\renewcommand{\slideid}{L3}
